% =============================================================================
% 02-entrevista.tex - PARTE 2, como perguntar. Frames 09 a 15.
% Frases LITERAIS na tela, nao parafrases.
% =============================================================================

% ----------------------------------------------------------------- frame 09 --
\begin{frame}{\ftit{\faCommentDots}{Perguntar n\~ao cria suicidalidade}}
\CRar{2.4}
\begin{caixaavaliar}{O receio e a evid\^encia}
O receio de que perguntar sobre suic\'idio possa induzir a ideia \'e comum e
n\~ao se sustenta. Revis\~ao da literatura publicada n\~ao encontrou, em
nenhum dos estudos revisados, aumento estatisticamente significativo de
idea\c{c}\~ao suicida entre os participantes que foram perguntados.
\end{caixaavaliar}

\CRvg
\begin{columns}[onlytextwidth, t]
\begin{column}{0.47\textwidth}
\begin{caixaproteger}{O que perguntar produz}
\begin{itemc}
\item nomeia o que o paciente n\~ao sabia como dizer
\item mostra que o assunto cabe nesta consulta
\item permite a pergunta seguinte
\end{itemc}
\end{caixaproteger}
\end{column}
\calha
\begin{column}{0.47\textwidth}
\begin{caixaagir}{O que n\~ao perguntar produz}
\begin{itemc}
\item um registro que diz apenas \fala{nega idea\c{c}\~ao}
\item uma alta que ningu\'em consegue justificar depois
\item a impress\~ao, no paciente, de que aquilo n\~ao se fala aqui
\end{itemc}
\end{caixaagir}
\end{column}
\end{columns}

\fonte{Dazzi, Gribble, Wessely e Fear, \emph{Psychol Med} 2014, revis\~ao em
formato de editorial. \seloevidencia}
\end{frame}

% ----------------------------------------------------------------- frame 10 --
\begin{frame}{\ftit{\faEye}{Reconhecer a pista}}
\CRar{2.6}
{\small Estas falas chegam ao pronto atendimento por outro motivo. Elas s\~ao a
abertura, e n\~ao a queixa.\par}

\CRvg
\begin{columns}[onlytextwidth, t]
\begin{column}{0.47\textwidth}
\begin{caixaavaliar}{O que o paciente diz}
\begin{itemc}
\item \fala{eu queria dormir e n\~ao acordar}
\item \fala{assim n\~ao d\'a mais}
\item \fala{eles ficariam melhor sem mim}
\item \fala{eu estou s\'o cansado, doutor}
\item \fala{n\~ao sei quanto tempo eu aguento}
\end{itemc}
\end{caixaavaliar}
\end{column}
\calha
\begin{column}{0.47\textwidth}
\begin{caixamudanca}{O que a fam\'ilia diz}
\begin{itemc}
\item \fala{ele deu as coisas dele}
\item \fala{ela ficou estranhamente calma}
\item \fala{ele se despediu de um jeito esquisito}
\item \fala{ela parou de atender o telefone}
\item \fala{ele acertou tudo o que estava pendente}
\end{itemc}
\end{caixamudanca}
\end{column}
\end{columns}

\mensagem{Toda pista reconhecida obriga \`a pergunta seguinte. Reconhecer e n\~ao
perguntar \'e pior do que n\~ao ter reconhecido.}
\end{frame}

% ----------------------------------------------------------------- frame 11 --
\begin{frame}{\ftit{\faComments}{A frase de transi\c{c}\~ao}}
\CRar{4.3}
{\footnotesize Esta \'e a frase. N\~ao \'e exemplo do que dizer, \'e o que se
diz.\par}

\CRvg
\begin{center}
\begin{tikzpicture}
  \node[musialBlue, text width=9.4cm, font=\footnotesize, inner ysep=6pt] {%
    \itshape\guillemotleft Pelo que voc\^e est\'a me contando, esse sofrimento
    est\'a bem pesado. Quando as pessoas chegam assim, \`as vezes pensam em
    morrer. Voc\^e tem tido esses pensamentos?\guillemotright};
\end{tikzpicture}
\end{center}

\CRvg
\begin{columns}[onlytextwidth, t]
\begin{column}{0.47\textwidth}
\begin{caixaproteger}{Por que funciona}
Nomeia o sofrimento antes de perguntar, universaliza sem banalizar, e faz uma
pergunta fechada que o paciente consegue responder com uma palavra.
\end{caixaproteger}
\end{column}
\calha
\begin{column}{0.47\textwidth}
\begin{caixaagir}{O que n\~ao dizer}
\fala{Voc\^e n\~ao est\'a pensando em fazer besteira, est\'a?}\par
\vspace{0.06cm}
A pergunta j\'a traz a resposta esperada.
\end{caixaagir}
\end{column}
\end{columns}
\end{frame}

% ----------------------------------------------------------------- frame 12 --
% Roteiro de perguntas. Uma das tres tabelas autorizadas na aula.
\begin{frame}{\ftit{\faSearch}{Caracterizar a idea\c{c}\~ao}}
\CRar{2.2}
\begin{CRtabela}
\begin{center}
\begin{tabular}{E P{4.05cm} P{4.05cm}}
\toprule
\rowcolor{campoazul!10}
\textbf{O que caracterizar} & \textbf{Como perguntar} & \textbf{O que muda} \\
\midrule
existe & \fala{voc\^e tem tido esses pensamentos?} & abre ou fecha toda a
sequ\^encia \\
frequ\^encia & \fala{quantas vezes, num dia como o de ontem?} & separa o
epis\'odico do cont\'inuo \\
dura\c{c}\~ao & \fala{quando vem, fica quanto tempo?} & minutos e horas pedem
condutas diferentes \\
intrusividade & \fala{ele vem sozinho ou voc\^e o chama?} & pensamento que
invade cansa e desgasta \\
capacidade de afastar & \fala{quando vem, voc\^e consegue empurrar? o que
ajuda?} & \'e aqui que mora o plano de seguran\c{c}a \\
\bottomrule
\end{tabular}
\end{center}
\end{CRtabela}

\mensagem{A \'ultima linha \'e a mais \'util das cinco, e \'e a que quase
ningu\'em pergunta.}
\end{frame}

% ----------------------------------------------------------------- frame 13 --
\begin{frame}{\ftit{\faRoute}{Inten\c{c}\~ao, plano e prepara\c{c}\~ao}}
\CRar{6}
\begin{center}
\begin{tikzpicture}[node distance=0.34cm]
  \node[crnode={campoazul}{2.55cm}, font=\scriptsize] (i) {\textbf{inten\c{c}\~ao}\\[1pt]
    \fala{voc\^e quer morrer ou quer que isso pare?}};
  \node[crnode={campoambar}{2.55cm}, font=\scriptsize, right=of i] (p) {\textbf{plano}\\[1pt]
    \fala{voc\^e chegou a pensar em como faria?}};
  \node[crnode={campoambar}{2.55cm}, font=\scriptsize, right=of p] (r) {\textbf{prepara\c{c}\~ao}\\[1pt]
    \fala{voc\^e chegou a fazer alguma coisa nessa dire\c{c}\~ao?}};
  \draw[crseta=campoazul] (i.east) -- (p.west);
  \draw[crseta=campoambar] (p.east) -- (r.west);
\end{tikzpicture}
\end{center}

\CRvg
\begin{columns}[onlytextwidth, t]
\begin{column}{0.47\textwidth}
\begin{caixaavaliar}{A pergunta que falta}
Depois de \fala{fez alguma coisa?}, vem \fala{mesmo que depois voc\^e tenha
desistido}. \'E essa segunda metade que autoriza o paciente a contar o ato
interrompido.
\end{caixaavaliar}
\end{column}
\calha
\begin{column}{0.47\textwidth}
\begin{caixaagir}{O que n\~ao entra na tela nem na pergunta}
N\~ao se pergunta detalhe de m\'etodo por curiosidade cl\'inica, e n\~ao se
compara letalidade diante do paciente. O que importa \'e a exist\^encia do
acesso e a possibilidade de restringi-lo.
\end{caixaagir}
\end{column}
\end{columns}
\end{frame}

% ----------------------------------------------------------------- frame 14 --
\begin{frame}{\ftit{\faClock}{Comportamento e temporalidade}}
\CRar{1.3}
\begin{caixamudanca}{A pergunta que organiza o tempo}
\fala{O que mudou nas \'ultimas horas? E nos \'ultimos dias?}
\end{caixamudanca}

\CRvg
\begin{center}
\begin{tikzpicture}[node distance=0.30cm]
  \node[crnode={campografite}{2.15cm}, font=\scriptsize] (v) {meses e anos\\[1pt]
    \scriptsize vulnerabilidade};
  \node[crnode={campoambar}{2.15cm}, font=\scriptsize, right=of v] (s) {\'ultimas semanas\\[1pt]
    \scriptsize deteriora\c{c}\~ao};
  \node[crnode={campoambar}{2.15cm}, font=\scriptsize, right=of s] (d) {\'ultimos dias\\[1pt]
    \scriptsize prepara\c{c}\~ao};
  \node[crnode={campobordo}{2.15cm}, font=\scriptsize, right=of d] (h) {\'ultimas horas\\[1pt]
    \scriptsize decis\~ao de hoje};
  \draw[musialThin] (v.east) -- (s.west);
  \draw[crseta=campoambar] (s.east) -- (d.west);
  \draw[crseta=campobordo] (d.east) -- (h.west);
\end{tikzpicture}
\end{center}

\CRvg
\begin{columns}[onlytextwidth, t]
\begin{column}{0.47\textwidth}
\begin{caixaavaliar}{O que perguntar do comportamento}
\begin{itemc}
\item o que fez, n\~ao s\'o o que pensou
\item o que interrompeu o ato, quando houve
\item o que estava acontecendo naquela hora
\end{itemc}
\end{caixaavaliar}
\end{column}
\calha
\begin{column}{0.47\textwidth}
\begin{caixamudanca}{Mudan\c{c}as que pesam}
\begin{itemc}
\item sono que encurtou de repente
\item retraimento que apertou nos \'ultimos dias
\item algo resolvido ou acertado sem explica\c{c}\~ao
\end{itemc}
\end{caixamudanca}
\end{column}
\end{columns}
\end{frame}

% ----------------------------------------------------------------- frame 15 --
\begin{frame}{\ftit{\faBolt}{O que modifica a seguran\c{c}a agora}}
\CRar{6}
\begin{center}
\begin{tikzpicture}[node distance=0.26cm and 0.30cm]
  \node[musialCenter, minimum size=1.95cm, font=\scriptsize\bfseries\sffamily] (c)
    {modificadores\\proximais};

  \node[crnode={campoambar}{1.85cm}, font=\scriptsize, above left=0.60cm and 0.75cm of c] (subs) {subst\^ancias};
  \node[crnode={campoambar}{1.85cm}, font=\scriptsize, above=0.60cm of c] (agit) {agita\c{c}\~ao};
  \node[crnode={campoambar}{1.85cm}, font=\scriptsize, above right=0.60cm and 0.75cm of c] (ins) {ins\^onia};
  \node[crnode={campobordo}{1.85cm}, font=\scriptsize, left=0.80cm of c] (psic) {psicose};
  \node[crnode={campobordo}{1.85cm}, font=\scriptsize, right=0.80cm of c] (imp) {impulsividade};
  \node[crnode={campoazul}{1.85cm}, font=\scriptsize, below left=0.60cm and 0.75cm of c] (perd) {perdas recentes};
  \node[crnode={campoverde}{1.85cm}, font=\scriptsize, below=0.60cm of c] (rede) {rede dispon\'ivel};
  \node[crnode={campoverde}{1.85cm}, font=\scriptsize, below right=0.60cm and 0.75cm of c] (colab) {capacidade de colaborar};

  \draw[musialThin] (subs.south east) -- (c.north west);
  \draw[musialThin] (agit.south) -- (c.north);
  \draw[musialThin] (ins.south west) -- (c.north east);
  \draw[musialThin] (psic.east) -- (c.west);
  \draw[musialThin] (imp.west) -- (c.east);
  \draw[musialThin] (perd.north east) -- (c.south west);
  \draw[musialThin] (rede.north) -- (c.south);
  \draw[musialThin] (colab.north west) -- (c.south east);
\end{tikzpicture}
\end{center}

\mensagem{Supervis\~ao dispon\'ivel e capacidade de colaborar s\~ao
modificadores, e n\~ao detalhe de execu\c{c}\~ao do plano.}
\end{frame}
